% Compile with latexmk (uses ~/.latexmkrc: LuaLaTeX, output in build/)
% Read-aloud teaching script for Week 10 (companion to beamer.tex, follows it
% frame by frame). Sessions 1 and 2 are held on the SAME DAY, back to back
% (2 x 90 min with a short break in between).
% Exercises are numbered by week (10.1, 10.2, ...); finance examples share the
% numbering. Each one lives in exercises/NN-slug.tex.
% The pictures and tables of the slides are copied into the script (after an
% "On the slide:" line), so that the script can be read without the deck.

\documentclass[11pt,a4paper,parskip=half]{scrartcl}
\usepackage{amsmath}
\usepackage[stixtwo]{fontsetup}
\usepackage[english]{babel}
\usepackage{booktabs}
\usepackage{enumitem}
\usepackage{needspace}
\usepackage{tikz}
\usepackage[most]{tcolorbox}
\usepackage[a4paper,margin=1in,marginparwidth=1.9cm,marginparsep=2mm]{geometry}
\usepackage[hidelinks]{hyperref}

\definecolor{navy}{RGB}{24,48,75}
\definecolor{teal}{RGB}{0,120,125}
\definecolor{gold}{RGB}{190,130,20}
\definecolor{plum}{RGB}{110,60,120}

\addtokomafont{disposition}{\color{navy}}
\setkomafont{section}{\Large\bfseries\color{navy}}

% Contact
\newcommand{\myemail}{tiborsandor2001@gmail.com}

% Exercise statement (what students see on the slide)
\newtcolorbox{task}[1]{enhanced,breakable,colback=navy!4,colframe=navy,
  fonttitle=\bfseries,title={#1}}
% Finance example
\newtcolorbox{fin}[1]{enhanced,breakable,colback=teal!5,colframe=teal,
  fonttitle=\bfseries,title={#1}}
% Teaching cue, not read aloud
\newtcolorbox{cue}{enhanced,breakable,colback=gold!8,colframe=gold,
  boxrule=.4pt,left=4pt,right=4pt,top=2pt,bottom=2pt,fontupper=\small\itshape}
% Background for the lecturer (finance theory), NOT read aloud
\newtcolorbox{bg}[1]{enhanced,breakable,colback=plum!4,colframe=plum,
  boxrule=.6pt,fonttitle=\bfseries\small,
  title={Background for you (not read aloud): #1},fontupper=\small,
  before upper={\setlength{\parskip}{4pt}}}
% Print version (script-print.tex defines \scriptprint): the background
% boxes are dropped entirely.
\ifdefined\scriptprint
  \RenewDocumentEnvironment{bg}{m +b}{}{}
\fi
\newcommand{\terms}{\par\textbf{\color{plum}Key terms.}\ }
\newcommand{\sq}[1]{\par\textbf{\color{plum}Student may ask:}\ \emph{#1}\par}
\newcommand{\ans}{\textbf{Answer:}\ }

\newcommand{\say}{\textbf{\color{teal}Say:}\ }
\newcommand{\board}{\textbf{\color{gold}Board:}\ }
\newcommand{\ask}{\textbf{\color{navy}Ask the class:}\ }
\newcommand{\slide}[1]{\marginpar{\raggedright\footnotesize\color{gray}slide #1}}
% What the slide shows (a picture or a table), copied into the script
\newcommand{\onslide}{\par\needspace{10\baselineskip}%
  \textbf{\color{gray}On the slide:}\ }

% The exercises (one file each in exercises/) and \scriptexercise
% Exercises of Week 10, shared by beamer.tex, the study guide, the script and
% exercises.tex. Every exercise lives in a file of its own,
% exercises/NN-slug.tex, numbered 10.NN (see exercises.lua for the docstring).
% The body of an exercise file:
%
%   \begin{problem} ... \end{problem}      the statement, used everywhere;
%                                          inside it, \begin{context} ...
%                                          \end{context} is background that
%                                          the slides leave out
%   \begin{slideproblem} ... \end{slideproblem}
%                                          (optional) a shorter statement
%                                          for the slides
%   \begin{slides} ... \end{slides}        one frame of the deck: \exstatement
%                                          places the statement; repeat the
%                                          environment for more frames
%   \begin{solution} ... \end{solution}    the detailed solution (study
%                                          guide, exercise PDFs, website)
%
% Commands for the documents (the argument is the slug: 01-license-plates.tex
% has the slug license-plates):
%
%   \exref{slug}                 the number, e.g. 10.1, for "Exercise~\exref{...}"
%   \extitle{slug}               the title from the docstring
%   \exerciseframe[opts]{slug}   beamer: the frames of the exercise
%   \studyexercise{slug}         study guide: heading, statement, solution
%   \scriptexercise{slug}        script: the statement box, as on the slide
%   \allexercises                exercises.tex: every exercise, in order; on
%                                the website statement and solution with
%                                markers where the pipeline splits the HTML
%                                (\collectionexercise), in the PDF plain:
%                                each statement, followed by its solution if
%                                \ifsolutions (exercises-solutions.tex)

\def\exweek{10}

\directlua{exercises = dofile("exercises/exercises.lua")
  exercises.scan("exercises", "\exweek")}

\ExplSyntaxOn

\msg_new:nnn { exercises } { unknown } { Unknown~exercise~`#1'. }

\cs_new_protected:Npn \__ex_missing:n #1
  { \msg_error:nnn { exercises } { unknown } {#1} ?? }
\cs_new:Npn \exref #1
  {
    \cs_if_exist:cTF { exnum@#1 }
      { \use:c { exnum@#1 } } { \__ex_missing:n {#1} }
  }
\cs_new:Npn \extitle #1
  {
    \cs_if_exist:cTF { extitle@#1 }
      { \use:c { extitle@#1 } } { \__ex_missing:n {#1} }
  }

% ---------------------------------------------------------------------------
% Loading an exercise: its parts are stored, nothing is typeset. The slides
% typeset the stored parts; the other documents read the part they need from
% the file again (\__ex_input_part:nn), as stored tables do not survive lwarp.

\tl_new:N \g__ex_current_tl
\tl_new:N \g__ex_problem_tl
\tl_new:N \g__ex_slideproblem_tl
\tl_new:N \g__ex_solution_tl
\seq_new:N \g__ex_slides_seq
\box_new:N \l__ex_scratch_box

\cs_new_protected:Npn \__ex_load:n #1
  {
    \cs_if_exist:cF { exfile@#1 } { \__ex_missing:n {#1} }
    \tl_gset:Nn \g__ex_current_tl {#1}
    \tl_gclear:N \g__ex_problem_tl
    \tl_gclear:N \g__ex_slideproblem_tl
    \tl_gclear:N \g__ex_solution_tl
    \seq_gclear:N \g__ex_slides_seq
    \vbox_set:Nn \l__ex_scratch_box
      {
        \DeclareDocumentEnvironment { problem } { +b }
          { \tl_gset:Nn \g__ex_problem_tl {##1} } { }
        \DeclareDocumentEnvironment { slideproblem } { +b }
          { \tl_gset:Nn \g__ex_slideproblem_tl {##1} } { }
        \DeclareDocumentEnvironment { slides } { +b }
          { \seq_gput_right:Nn \g__ex_slides_seq {##1} } { }
        \DeclareDocumentEnvironment { solution } { +b }
          { \tl_gset:Nn \g__ex_solution_tl {##1} } { }
        \exp_args:Nv \input { exfile@#1 }
      }
  }

% Typesets one part (problem, slideproblem, solution) of an exercise straight
% from its file; the other parts are skipped.
\cs_new_protected:Npn \__ex_input_part:nn #1#2
  {
    \group_begin:
      \clist_map_inline:nn { problem, slideproblem, slides, solution }
        {
          \str_if_eq:nnTF {##1} {#2}
            { \DeclareDocumentEnvironment {##1} { } { } { } }
            { \DeclareDocumentEnvironment {##1} { +b } { } { } }
        }
      \exp_args:Nv \input { exfile@#1 }
    \group_end:
  }

% Exercise 10.1, or Exercise 10.2 · Finance
\cs_new:Npn \__ex_heading:
  {
    Exercise~\exref { \g__ex_current_tl }
    \cs_if_exist:cT { exfinance@ \g__ex_current_tl } { ~·~Finance }
  }
\prg_new_conditional:Npnn \__ex_if_finance: { T, F, TF }
  {
    \cs_if_exist:cTF { exfinance@ \g__ex_current_tl }
      { \prg_return_true: } { \prg_return_false: }
  }

% ---------------------------------------------------------------------------
% Beamer

% The finance context is printed in the study guide, the exercise PDFs and on
% the website; the
% slides and the script (which shows what the slide shows) leave it out.
\bool_new:N \l__ex_context_bool
\bool_set_true:N \l__ex_context_bool
\DeclareDocumentEnvironment { context } { }
  {
    \bool_if:NTF \l__ex_context_bool
      { \par \textit { Finance~context. } ~ }
      { \__ex_gobble_context:w }
  }
  { \bool_if:NT \l__ex_context_bool { \par } }
\cs_new_protected:Npn \__ex_gobble_context:w #1 \end #2
  { \str_if_eq:nnTF {#2} { context } { \end { context } } { \__ex_gobble_context:w } }

\IfClassLoadedTF { beamer }
  {
    \bool_set_false:N \l__ex_context_bool

    \cs_new:Npn \__ex_block: { \__ex_if_finance:TF { exampleblock } { block } }
    \NewDocumentCommand \exstatement { }
      {
        \begin { \__ex_block: } { \__ex_heading: }
          \tl_if_empty:NTF \g__ex_slideproblem_tl
            { \g__ex_problem_tl } { \g__ex_slideproblem_tl }
        \end { \__ex_block: }
      }

    \tl_new:N \l__ex_frames_tl
    \NewDocumentCommand \exerciseframe { O{t} m }
      {
        \__ex_load:n {#2}
        \tl_clear:N \l__ex_frames_tl
        \seq_map_indexed_inline:Nn \g__ex_slides_seq
          {
            \tl_put_right:Nn \l__ex_frames_tl { \begin{frame}[#1] }
            \tl_put_right:Ne \l__ex_frames_tl
              {
                {
                  \exp_not:N \extitle {#2}
                  \int_compare:nNnT {##1} > { 1 } { \c_space_tl (cont.) }
                }
              }
            \tl_put_right:Nn \l__ex_frames_tl { ##2 \end{frame} }
          }
        \tl_use:N \l__ex_frames_tl
      }
  }
  {
    \cs_new:Npn \__ex_box: { \__ex_if_finance:TF { fin } { task } }
    \cs_new_protected:Npn \__ex_statement_box:n #1
      {
        \begin { \__ex_box: } { \__ex_heading: }
          \__ex_input_part:nn { \g__ex_current_tl } {#1}
        \end { \__ex_box: }
      }

    \NewDocumentCommand \studyexercise { m }
      {
        \__ex_load:n {#1}
        \subsection* { Exercise~\exref{#1}~---~\extitle{#1} }
        \addcontentsline { toc } { subsection }
          { Exercise~\exref{#1}~---~\extitle{#1} }
        \__ex_statement_box:n { problem }
        \soltitle \__ex_input_part:nn {#1} { solution }
      }

    \NewDocumentCommand \scriptexercise { m }
      {
        \__ex_load:n {#1}
        \group_begin:
          \bool_set_false:N \l__ex_context_bool
          \tl_if_empty:NTF \g__ex_slideproblem_tl
            { \__ex_statement_box:n { problem } }
            { \__ex_statement_box:n { slideproblem } }
        \group_end:
      }

    % The website (lwarp): statement and solution of each exercise, between
    % markers where the pipeline splits the HTML.
    \NewDocumentCommand \collectionexercise { m }
      {
        \__ex_load:n {#1}
        \par [[exercise:#1:problem]] \par
        \__ex_input_part:nn {#1} { problem }
        \par [[exercise:#1:solution]] \par
        \__ex_input_part:nn {#1} { solution }
        \par [[exercise:#1:end]] \par
      }
    % The PDF: plain, each statement, followed by its solution in the
    % version with solutions (\ifsolutions). The statement is set off like a
    % quote in Markdown, by a bar on the left (quote=<color>, also used by
    % exercises.tex) in the color of its box in the study guide: navy, or teal
    % for a finance exercise.
    \tcbset
      {
        quote/.style =
          {
            enhanced, breakable, frame~hidden, interior~hidden,
            sharp~corners, boxrule = 0pt, borderline~west = {2pt}{0pt}{#1},
            left = 8pt, right = 0pt, top = 1pt, bottom = 1pt,
            before~skip = 6pt, after~skip = 6pt,
            % the document's paragraph spacing (parskip=half)
            before~upper = { \setlength { \parskip } { .5\baselineskip } }
          }
      }
    \NewDocumentCommand \plainexercise { m }
      {
        \__ex_load:n {#1}
        \subsection* { Exercise~\exref{#1}~---~\extitle{#1} }
        \addcontentsline { toc } { subsection }
          { Exercise~\exref{#1}~---~\extitle{#1} }
        \__ex_if_finance:TF
          { \begin{tcolorbox}[quote=teal] }
          { \begin{tcolorbox}[quote=navy] }
          \__ex_input_part:nn {#1} { problem }
        \end{tcolorbox}
        \legacy_if:nT { solutions }
          {
            \soltitle
            \__ex_input_part:nn {#1} { solution }
          }
      }
    \NewDocumentCommand \allexercises { }
      {
        \IfPackageLoadedTF { lwarp }
          { \directlua { exercises.each("collectionexercise") } }
          { \directlua { exercises.each("plainexercise") } }
      }
  }

\ExplSyntaxOff


% PDF properties and the date on the title page
% The properties of every PDF of Week 10: author (with email address), title,
% subject, keywords and the date the document last changed, in the PDF
% (hyperref) and on its title page (\date), and the header and footer of the
% printed documents (below). Each document loads this file after hyperref and
% then calls \documentinfo{<what it is>}, e.g. \documentinfo{Slides}.
%
% The date: the website build (scripts/build-content.mjs) puts
% \def\lastmodified{YYYY-MM-DD} in front of the document, the date of the last
% commit that changed its sources. Compiled locally, it is today.

\newcommand{\courseinfo}{Mathematics for Finance}
\newcommand{\weeklabel}{Week 10}
\newcommand{\weekinfo}{\weeklabel: Continuous random variables, joint distributions, independence}
\newcommand{\authorname}{Tibor SÁNDOR}
\providecommand{\myemail}{tiborsandor2001@gmail.com}

\makeatletter
\ifdefined\lastmodified\else
	\edef\lastmodified{\the\year-\two@digits{\the\month}-\two@digits{\the\day}}
\fi
\def\di@split#1-#2-#3\@nil{%
	\def\di@year{#1}\def\di@month{#2}\def\di@day{#3}}
\expandafter\di@split\lastmodified\@nil
% November 11, 2026
\edef\lastmodifiedtext{%
	\ifcase\numexpr\di@month\relax\or January\or February\or March\or April\or
		May\or June\or July\or August\or September\or October\or November\or
		December\fi\space\number\di@day, \di@year}
% D:20261111000000, the date format of PDF metadata
\edef\di@pdfdate{D:\di@year\di@month\di@day000000}

\newcommand{\documentinfo}[1]{%
	\hypersetup{
		pdftitle={\courseinfo, \weekinfo\space -- #1},
		pdfauthor={\authorname\space<\myemail>},
		pdfsubject={\courseinfo, \weekinfo},
		pdfkeywords={probability, random variables, mathematics for finance,
				Corvinus University of Budapest},
		pdfcreationdate={\di@pdfdate},
		pdfmoddate={\di@pdfdate},
		pdflang={en-US},
	}%
	\date{Last updated: \lastmodifiedtext}%
	\gdef\documentkind{#1}%
}

% Header and footer of the printed documents (KOMA-Script, scrlayer-scrpage;
% the slides have beamer's footline, the website neither): the course, the
% week and the document on the left of the header (the week's topic is on the
% title page), the university on the right; in the footer the author and the
% page as <current>/<total>. The title page gets only the footer.
\newcommand{\documentkind}{}
\IfClassLoadedTF{beamer}{}{\IfPackageLoadedTF{lwarp}{}{%
	\usepackage{lastpage}%
	\usepackage[headsepline=1.5pt,footsepline=1.5pt,plainfootsepline]
		{scrlayer-scrpage}%
	\clearpairofpagestyles
	\setkomafont{pageheadfoot}{\small\sffamily\color{navy}}%
	\setkomafont{pagenumber}{}%
	\setkomafont{headsepline}{\color{navy!35}}%
	\setkomafont{footsepline}{\color{navy!35}}%
	\ihead{\courseinfo\ \textperiodcentered\ \weeklabel\ \textperiodcentered\ 
		\documentkind}%
	\ohead{Corvinus University of Budapest}%
	\ifoot*{\authorname\ \textperiodcentered\ \myemail}%
	\ofoot*{\pagemark/\pageref*{LastPage}}%
	\pagestyle{scrheadings}%
	% Extra stretch for breaking lines with long inline formulas.
	\setlength{\emergencystretch}{1.5em}%
}}
\makeatother

\documentinfo{Teaching script\ifdefined\scriptprint\space(print version)\fi}

\title{Mathematics for Finance}
\subtitle{Week 10 --- Continuous random variables, joint distributions,
  independence\\[2pt]
  \normalsize Teaching script for two 90-minute practical classes held on the
  same day, back to back\ifdefined\scriptprint\\[2pt]
  (print version, without the background boxes)\fi}
\author{Tibor SÁNDOR\\[2pt]
  \normalsize\href{mailto:\myemail}{\texttt{\myemail}}}

\begin{document}
\maketitle

\begin{cue}
	How to use this script. Plain text is meant to be read out loud or
	paraphrased. \textbf{Say} marks key sentences; \textbf{Ask} marks questions
	to put to the class (wait for answers before revealing the solution);
	\textbf{Board} marks what to write. Gold boxes like this one are short notes
	for you only.\ifdefined\scriptprint\else\ \textbf{Purple ``Background for
			you'' boxes} explain the finance behind the examples; they are
		\emph{not} meant to be read out.\fi\ The script follows the deck frame by
	frame; the slide numbers in the margin are the frame numbers in the
	footline of the deck. Every picture and table of the slides is copied
	into the script after an \textbf{On the slide} line, and the formulas of
	each theory slide are displayed, so you can teach from the script alone.
	Every new idea gets one tiny example; before each exercise the students
	try it in pairs, then you reveal the solution step by step. Only the
	exercises of the slides are in this script.
\end{cue}

\subsection*{Timing overview}

Both sessions are on the same day: Session 1 (90 min), a short break,
Session 2 (90 min). The plan leaves a generous reserve: 20 minutes in
Session 1 and 25 in Session 2, for questions and for students working
alone.

\begin{center}
	\small
	\begin{tabular}{@{}llrr@{}}
		\toprule
		Slides & Session 1 & min & cum. \\\midrule
		1--4   & Welcome, recap, plan, from discrete to continuous & 6  & 6  \\
		5      & Continuous random variables and their density     & 5  & 11 \\
		6--7   & Exercise~\exref{settlement-time}                  & 10 & 21 \\
		8      & The distribution function                         & 5  & 26 \\
		9--10  & Exercise~\exref{loss-given-default}               & 10 & 36 \\
		11     & Uniform and exponential                           & 5  & 41 \\
		12--13 & Exercise~\exref{order-arrivals}                   & 10 & 51 \\
		14     & The distribution of a function of a random variable & 6 & 57 \\
		15--16 & Exercise~\exref{square-of-uniform}                & 10 & 67 \\
		17     & Summary                                           & 3  & 70 \\
		       & Reserve (questions, students working alone)       & 20 & 90 \\\midrule
		       & \textbf{Total}                                    &    & \textbf{90} \\\bottomrule
	\end{tabular}
\end{center}

\begin{center}
	\small
	\begin{tabular}{@{}llrr@{}}
		\toprule
		Slides & Session 2 & min & cum. \\\midrule
		18--19 & Session 2 title, recap of Session 1               & 4  & 4  \\
		20     & Joint distribution: the discrete case             & 6  & 10 \\
		21--22 & Exercise~\exref{two-stocks}                       & 10 & 20 \\
		23     & Joint distribution: the continuous case           & 6  & 26 \\
		24--25 & Exercise~\exref{joint-density-square}             & 12 & 38 \\
		26     & Independent random variables                      & 6  & 44 \\
		27--28 & Exercise~\exref{two-stocks-independence}          & 8  & 52 \\
		29--30 & Exercise~\exref{two-default-times}                & 10 & 62 \\
		31     & Summary                                           & 3  & 65 \\
		       & Reserve (questions, students working alone)       & 25 & 90 \\\midrule
		       & \textbf{Total}                                    &    & \textbf{90} \\\bottomrule
	\end{tabular}
\end{center}

\subsection*{If you are short on time}

\begin{cue}
	Skip \emph{parts} of exercises, never a whole topic, in this order. Every
	skipped part has a full solution in the exercises PDF: tell the students
	so. Gold \emph{Clock check} boxes in the script say when to decide.

	\textbf{Session 1:}
	\begin{enumerate}
		\item Part (d) of Exercise~\exref{loss-given-default}, the median
		      (slide 10, about 2 min). Decide at slide 8.
		\item Part (c) of Exercise~\exref{square-of-uniform} (slide 16): just
		      give the two numbers, $0.5$ and $0.1$. Decide at slide 14.
	\end{enumerate}

	\textbf{Session 2:}
	\begin{enumerate}
		\item Part (d) of Exercise~\exref{two-default-times} (slide 30, about
		      2 min). Decide at slide 23.
		\item Part (d) of Exercise~\exref{joint-density-square} (slide 25):
		      state the symmetry argument in one sentence. Decide at slide 23.
		\item Part (c) of Exercise~\exref{two-stocks-independence} (slide 28):
		      give $0.40$ against $0.60$ and the one-sentence lesson. Decide at
		      slide 26.
	\end{enumerate}

	If you are behind everywhere, shorten the try-first pauses: one or two
	minutes instead of three.
\end{cue}

\tableofcontents

%======================================================================
\section{Session 1: Continuous random variables}
%======================================================================

\subsection{Welcome, recap, plan, from discrete to continuous (6 min)}\slide{1--4}

\begin{cue}
	Slide 1 (title) while they settle down.
\end{cue}

Good morning, everyone, and welcome back.

\say Last week we met probability, conditional probability, independence,
and discrete random variables: numbers that depend on chance and take
values like $0,1,2,\dots$. A discrete random variable is described by its
pmf, $p(x)=P(X=x)$, and the probabilities add up to 1.

\begin{cue}
	Slide 2 (Plan for today).
\end{cue}

\say Today, Session 1: random variables that can take \emph{any} value in
an interval, described by a density; the distribution function; two
important examples, the uniform and the exponential distribution; and how
to find the distribution of a function of a random variable. Session 2:
two random variables at the same time, and when they are independent.

Why do we need this in finance? Waiting times, losses and returns can
take any value in an interval, and the risk of a portfolio depends on how
its parts move together.

\begin{cue}
	As last week: before every exercise they get a few minutes to try in
	pairs. Slide 3 (Session 1 title): click through to slide 4.
\end{cue}

\say Think of the time until the next order arrives for a stock. It could
be $0.5$ minutes, $0.51$, $0.5132$\ldots{} There are far too many values to
give each one a positive probability: the sum would explode. So for a
continuous quantity, $P(X=a)=0$ for every single number $a$.

Instead, we ask for the probability of an \emph{interval},
$P(a\le X\le b)$, and we describe it with a curve, the density. The
probability is the \emph{area under the curve} between $a$ and $b$.

\onslide a bell-shaped density $f(x)$; the shaded area between $a$ and $b$
is the probability $P(a\le X\le b)$.
\begin{center}
	\begin{tikzpicture}[scale=.9]
		\fill[teal!25] (1.2,0) -- plot[domain=1.2:2.6,samples=40]
			(\x,{1.6*exp(-0.5*(\x-2)^2)}) -- (2.6,0) -- cycle;
		\draw[->] (-.2,0) -- (4.6,0) node[right] {$x$};
		\draw[thick,navy] plot[domain=0:4.3,samples=60]
			(\x,{1.6*exp(-0.5*(\x-2)^2)});
		\node[navy] at (3.9,1.2) {$f(x)$};
		\node[below] at (1.2,0) {\small $a$};
		\node[below] at (2.6,0) {\small $b$};
		\node at (1.9,.55) {\small area};
	\end{tikzpicture}
\end{center}

\say Look at the picture: the curve itself is not a probability. The
probability lives in the shaded \emph{area}. Where the curve is high,
intervals collect a lot of area, so values there are likely; where it is
low, values are unlikely. Think of a histogram of many observed waiting
times with narrower and narrower bars: in the limit, the tops of the bars
become this smooth curve.

\subsection{Continuous random variables and their density (5 min)}\slide{5}

\onslide the definition, the two properties of a density, and two
remarks: $P(X=a)=\int_a^a f=0$, so $P(a\le X\le b)=P(a<X<b)$; and $f(x)$ is
\emph{not} a probability, it can be bigger than 1. The formulas are
displayed below.

\say $X$ is continuous if there is a function $f$, its density, such that
\[ P(a\le X\le b)=\int_a^b f(x)\,dx. \]
A density has two properties: $f(x)\ge0$ everywhere, and the total area
is 1:
\[ \text{(1) } f(x)\ge0 \text{ for every } x, \qquad
	\text{(2) } \int_{-\infty}^{\infty}f(x)\,dx=1. \]

Why exactly these two? Probabilities are never negative, so the area over
any interval must be nonnegative, and that forces $f\ge0$. And $X$ takes
\emph{some} value with certainty, so the area over the whole line, the
probability of everything, must be 1. These two conditions are all we ever
check when someone hands us a function and says ``this is a density''.

Tiny example: $f(x)=3x^2$ on $[0,1]$, zero elsewhere. Total area:
$\int_0^1 3x^2\,dx=\big[x^3\big]_0^1=1$. And
$P(X\le\tfrac12)=\int_0^{1/2}3x^2\,dx=\tfrac18$. So the left half of the
interval carries only one eighth of the probability, the right half seven
eighths: the density grows, and large values are more likely.

\say Two warnings. First, a single point has no area, so $P(X=a)=0$, and
for a continuous $X$ it does not matter whether we write $<$ or $\le$.
Second, $f(x)$ itself is not a probability. In the example $f(1)=3$, bigger
than 1. Only areas are probabilities.

What does $f(1)=3$ mean then? It is a \emph{rate}: probability per unit
length. For a short interval of length $h$ near 1, the probability is about
$f(1)\cdot h$. With $h=0.01$:
$P(0.99\le X\le1)=1-0.99^3\approx0.0297$, very close to $3\cdot0.01=0.03$.
A density of 3 over a short stretch gives a small probability.

\subsection{Exercise~\exref{settlement-time} --- Settling a payment (10 min)}\slide{6--7}

\scriptexercise{settlement-time}

\say Try (a) and (b) in pairs, three minutes. Hint: what must the total
area be?

\begin{cue}
	While they work, walk around. Typical problem: they integrate
	$c\,x(2-x)$ but do not know what to set it equal to. Point them back to
	the second property of a density. After 3 minutes, click: the two
	properties and ``probability $=$ area''.
\end{cue}

\say Let us start with the rule we use in every part: a density is
nonnegative, its total area is 1, and every probability is an area under
it.

(a) Click. \ask What condition fixes $c$? \emph{(The total area is 1.)}
Outside $[0,2]$ the density is zero, so the total area is the integral
over $[0,2]$. We integrate first without $c$; multiplying out,
$x(2-x)=2x-x^2$:
\[ \int_0^2 x(2-x)\,dx=\Big[x^2-\frac{x^3}{3}\Big]_0^2=4-\frac83=\frac43. \]
Click. So $c\cdot\tfrac43=1$ and $c=\tfrac34$.

Quick check of the first property: on $[0,2]$ both $x$ and $2-x$ are
nonnegative, so $f(x)=\tfrac34x(2-x)\ge0$. The density is a parabola arch
with its top $f(1)=\tfrac34$ at $x=1$: settling takes about one hour most
often.

\begin{cue}
	Slide 7.
\end{cue}

\onslide the density $f(x)=\tfrac34x(2-x)$ on $[0,2]$, a parabola arch
with height $0.75$ at $x=1$, a dashed gold line at $x=1$; from the second
click, the area between $0.5$ and $1.5$ is shaded, labeled $0.6875$.
\begin{center}
	\begin{tikzpicture}[xscale=2.4,yscale=3.6]
		\fill[teal!30] (.5,0) -- plot[domain=.5:1.5,samples=40]
			(\x,{.75*\x*(2-\x)}) -- (1.5,0) -- cycle;
		\node[teal!70!black] at (1,.28) {\small $0.6875$};
		\draw[->] (-.2,0) -- (2.35,0) node[right] {\small $x$};
		\draw[->] (0,-.05) -- (0,.95) node[above] {\small $f(x)$};
		\draw[navy,thick] plot[domain=0:2,samples=50] (\x,{.75*\x*(2-\x)});
		\foreach \x in {0.5,1,1.5,2}
		\draw (\x,.02) -- (\x,-.02) node[below] {\scriptsize $\x$};
		\draw (.04,.75) -- (-.04,.75) node[left] {\scriptsize $0.75$};
		\draw[gold,dashed] (1,0) -- (1,.75);
	\end{tikzpicture}
\end{center}

(b) \ask Look at the graph. Do we need to integrate? \emph{(No: the graph
	is symmetric around $x=1$, so the area splits in half.)}
$P(X>1)=\tfrac12$.

\say The dashed line at $x=1$ cuts the arch into two mirror images: the
area to the right equals the area to the left, and together they are 1.
In formulas, $f(1+t)=\tfrac34(1+t)(1-t)=f(1-t)$. Whenever you see a
symmetric density, use it: it saves an integral.

\ask Before we compute (c): will $P(0.5<X<1.5)$ be more or less than
$\tfrac12$? \emph{(More: the interval has half the length of $[0,2]$, but
it sits in the middle, where the density is highest.)}

(c) Click. Now we integrate $\tfrac34(2x-x^2)$ from $0.5$ to $1.5$. The
shaded area in the picture:
\[ \frac34\Big[x^2-\frac{x^3}{3}\Big]_{0.5}^{1.5}
	=\frac34\Big(\frac98-\frac5{24}\Big)=\frac34\cdot\frac{22}{24}
	=\frac{11}{16}=0.6875. \]
The two values of the antiderivative: at $1.5$ it is
$2.25-1.125=\tfrac98$, at $0.5$ it is $0.25-\tfrac{0.125}{3}=\tfrac5{24}$.
So with probability $0.6875$ the payment is settled between half an hour
and one and a half hours after it arrives; our guess ``more than half''
was right.

(d) Click. \ask $P(X=1)$? \emph{(Zero: a single point has no area.)}
So $P(X<1)=P(X\le1)$.

\say Zero probability does not mean that settling in exactly one hour is
impossible. It means that one exact value carries no area. And because of
that, in (c) it makes no difference whether we write $0.5<X<1.5$ or
$0.5\le X\le1.5$.

\begin{cue}
	Common mistake: reading $f(1)=0.75$ as a probability. The value of a
	density is not a probability; only areas are. Another one in (c):
	forgetting the factor $\tfrac34$ and getting $\tfrac{11}{12}$.
\end{cue}

\subsection{The distribution function (5 min)}\slide{8}

\begin{cue}
	Clock check: this slide is planned for minute 21. If you start it after
	about minute 26, skip part (d) of Exercise~\exref{loss-given-default}.
\end{cue}

\say So far every probability needed an integral. The distribution
function stores all these integrals in one function, so that afterward
we only subtract.

\onslide the definitions, and on the right a picture.
\[ F(x)=P(X\le x),\qquad F(x)=\int_{-\infty}^{x}f(t)\,dt,\qquad f(x)=F'(x), \]
\[ P(a<X\le b)=F(b)-F(a),\qquad P(X>a)=1-F(a). \]
The picture: an S-shaped curve $F(x)$ that starts at 0 and rises to the
dashed level 1; two points $a<b$ on the $x$-axis, dotted lines up to the
curve, and a teal arrow on the right showing the height difference
$F(b)-F(a)$.
\begin{center}
	\begin{tikzpicture}[scale=.9]
		\draw[->] (-.3,0) -- (4.2,0) node[right] {$x$};
		\draw[->] (0,-.1) -- (0,2.3);
		\draw[dashed,gray] (0,2) node[left,black] {\small 1} -- (4.1,2);
		\draw[thick,navy] (-.3,0) -- (.5,0)
			plot[domain=.5:3.5,samples=40]
			(\x,{2*(3*((\x-.5)/3)^2-2*((\x-.5)/3)^3)}) -- (4.1,2);
		\node[navy] at (.9,1.4) {$F(x)$};
		\draw[dotted] (1.6,0) node[below] {\small $a$} -- (1.6,.52)
			-- (4.1,.52);
		\draw[dotted] (2.8,0) node[below] {\small $b$} -- (2.8,1.55)
			-- (4.1,1.55);
		\draw[<->,teal,thick] (4.1,.52) -- (4.1,1.55);
		\node[teal,right] at (4.1,1.03) {\small $F(b)-F(a)$};
	\end{tikzpicture}
\end{center}

\say For any random variable, discrete or continuous, the distribution
function is $F(x)=P(X\le x)$. For a continuous $X$ it is the area to the left of $x$:
$F(x)=\int_{-\infty}^x f(t)\,dt$, and by the fundamental theorem of
calculus $f(x)=F'(x)$. Integrate the density to get $F$; differentiate $F$
to get the density.

\board $P(a<X\le b)=F(b)-F(a)$, \quad $P(X>a)=1-F(a)$.

Why the first rule? The event $X\le b$ splits into two disjoint pieces,
$X\le a$ and $a<X\le b$. So $F(b)=F(a)+P(a<X\le b)$, and we subtract. In the
picture: the probability of the interval is the height of the teal arrow,
the rise of $F$ between $a$ and $b$. The second rule is the complement
rule: ``greater than $a$'' is everything except ``at most $a$''.

Tiny example: for $f(x)=3x^2$ on $[0,1]$, $F(x)=x^3$ on $[0,1]$, so
$P(X>\tfrac12)=1-\tfrac18=\tfrac78$. And
$P(\tfrac14<X\le\tfrac12)=\tfrac18-\tfrac1{64}=\tfrac7{64}$, with no integral
at all.

\say $F$ grows from 0 to 1. For a continuous $X$ it has no jumps; last
week, for a discrete $X$, it was a step function.

\ask Where is $F$ steepest in the picture? \emph{(In the middle, where
the density is largest: the slope of $F$ is the density.)}

\subsection{Exercise~\exref{loss-given-default} --- Loss given default, finance (10 min)}\slide{9--10}

\scriptexercise{loss-given-default}

\say The loss given default is the fraction of the loan that the bank
loses if the borrower defaults: 0 means nothing lost, 1 means everything.
Try (a)--(c), three minutes.

\begin{cue}
	If a pair is stuck, ask them: ``How are $f$ and $F$ related?''
\end{cue}

(a) Click. The density is the derivative of $F$: $f(x)=2x$ on $[0,1]$, $0$
otherwise. Outside $[0,1]$, $F$ is constant (0 on the left, 1 on the
right), and the derivative of a constant is 0. Click. Check:
$\int_0^1 2x\,dx=\big[x^2\big]_0^1=1$. The density grows, so large losses
are more likely than small ones.

\onslide the two functions in one picture: the density $f(x)=2x$ (navy), a
straight line from 0 up to 2 at $x=1$, then 0; and the distribution
function $F(x)=x^2$ (gold), a curve from 0 up to 1, then flat at 1.
\begin{center}
	\begin{tikzpicture}[xscale=2.2,yscale=1.2]
		\draw[->] (-.1,0) -- (1.4,0) node[right] {\small $x$};
		\draw[->] (0,-.1) -- (0,2.3) node[above] {\small $y$};
		\draw[gold,thick] plot[domain=0:1,samples=50] (\x,{\x*\x});
		\draw[navy,thick] (0,0) -- (1,2);
		\draw[navy,thick] (1,0) -- (1.3,0);
		\draw[navy,dashed] (1,0) -- (1,2);
		\draw[gold,thick] (1,1) -- (1.3,1);
		\draw (1,.04) -- (1,-.04) node[below] {\scriptsize $1$};
		\draw (.04,1) -- (-.04,1) node[left] {\scriptsize $1$};
		\draw (.04,2) -- (-.04,2) node[left] {\scriptsize $2$};
		\node[navy,left] at (.95,2.08) {\small $f(x)=2x$};
		\node[gold!80!black,right] at (1.02,1.25) {\small $F(x)=x^2$};
	\end{tikzpicture}
\end{center}

\say In the picture the density, the straight navy line, is twice as
high at 1 as at $0.5$; the gold curve $F$ is flat at the start and steep at
the end. Steep $F$ means high density: most of the probability sits near
1. In this model, when this borrower defaults, the bank typically loses a
large part of the loan.

\begin{cue}
	Slide 10, one part per click.
\end{cue}

\onslide the density $f(x)=2x$ again; the area under it between $0.5$ and
1 is shaded and labeled $0.75$ (``shaded area $=P(L>0.5)$'').
\begin{center}
	\begin{tikzpicture}[scale=1.7]
		\fill[teal!30] (.5,0) -- (.5,1) -- (1,2) -- (1,0) -- cycle;
		\draw[->] (-.1,0) -- (1.3,0) node[right] {\small $x$};
		\draw[->] (0,-.1) -- (0,2.25) node[above] {\small $f(x)$};
		\draw[navy,thick] (0,0) -- (1,2);
		\draw[navy,dashed] (1,0) -- (1,2);
		\draw (.5,.04) -- (.5,-.04) node[below] {\scriptsize $0.5$};
		\draw (1,.04) -- (1,-.04) node[below] {\scriptsize $1$};
		\draw (.04,2) -- (-.04,2) node[left] {\scriptsize $2$};
		\node[teal!60!black] at (.78,.55) {\small $0.75$};
	\end{tikzpicture}
\end{center}

\ask Before the click: a student says $P(L>0.5)=0.5$, ``because $0.5$ is
half of the interval''. Right or wrong? \emph{(Wrong: that would be true
for a uniform $L$; here the right half carries more area.)}

(b) Complement:
$P(L>0.5)=1-F(0.5)=1-0.5^2=0.75$.

\say The shaded region in the picture is a trapezoid with parallel sides
$f(0.5)=1$ and $f(1)=2$ and width $0.5$: area
$\tfrac{1+2}{2}\cdot0.5=0.75$. Geometry and $F$ give the same number.

(c) No integral needed: $P(0.2<L\le0.6)=F(0.6)-F(0.2)=0.36-0.04=0.32$.

(d) \ask Which equation gives the median? \emph{($F(m)=\tfrac12$.)}
So $m^2=\tfrac12$, $m=\sqrt{0.5}\approx0.707$: half of the defaults lose
more than 70.7\% of the loan.

\say We take the positive root because $m$ must lie between 0 and 1.

\begin{cue}
	Common mistakes: (b) as $1-0.5=0.5$, as if $L$ were uniform; (c) as
	$F(0.6-0.2)=F(0.4)=0.16$. An interval is the \emph{difference} of $F$
	at the endpoints.
\end{cue}

\begin{bg}{loss given default}
	\emph{Loss given default (LGD)}: the share of the exposure a lender
	loses when the borrower defaults, after collateral and recoveries.
	$\text{LGD}=1-\text{recovery rate}$. Typical values: around
	$0.2$--$0.45$ for secured or senior loans, $0.6$--$0.8$ for unsecured or
	subordinated debt. Expected loss $=\text{PD}\times\text{LGD}\times
		\text{EAD}$ (exposure at default); under the Basel rules banks
	estimate all three. LGD is a random variable between 0 and 1, so it is
	modeled with a density on $[0,1]$ (in practice often a beta
	distribution; $F(x)=x^2$ here is a simple special case).

	\terms loss given default (LGD); recovery rate; exposure at default
	(EAD); expected loss.

	\sq{Why would large losses be more likely?}
	\ans For unsecured loans, recoveries are often small: the bank gets
	back little. Then LGD piles up near 1.
\end{bg}

\subsection{Two examples: uniform and exponential (5 min)}\slide{11}

\say Two distributions we will use again and again.

\onslide two blocks with a picture under each.
\[ \text{Uniform on }[\alpha,\beta]:\quad f(x)=\frac{1}{\beta-\alpha}
	\text{ on }[\alpha,\beta],\ 0\text{ outside};\qquad
	\text{probability}=\frac{\text{length}}{\beta-\alpha}. \]
\[ \text{Exponential with rate }\lambda>0:\quad
	f(x)=\lambda e^{-\lambda x}\ (x\ge0),\qquad F(x)=1-e^{-\lambda x},\qquad
	P(X>x)=e^{-\lambda x}. \]
Left picture: the uniform density, a flat line between $\alpha$ and
$\beta$, with a shaded rectangle over a small interval inside. Right
picture: the exponential density, a decreasing curve; the area to the
right of a point $x$ is shaded and labeled $e^{-\lambda x}$.
\begin{center}
	\begin{tikzpicture}[scale=.85]
		\draw[->] (-.3,0) -- (4.2,0) node[right] {$x$};
		\fill[teal!25] (1.6,0) rectangle (2.6,1.2);
		\draw[thick,navy] (-.3,0) -- (.6,0) (.6,1.2) -- (3.4,1.2)
			(3.4,0) -- (4,0);
		\draw[dotted] (.6,0) node[below] {\small $\alpha$} -- (.6,1.2);
		\draw[dotted] (3.4,0) node[below] {\small $\beta$} -- (3.4,1.2);
	\end{tikzpicture}
	\hspace{1.5cm}
	\begin{tikzpicture}[scale=.85]
		\draw[->] (-.3,0) -- (4.2,0) node[right] {$x$};
		\fill[teal!25] (1.5,0) -- plot[domain=1.5:4,samples=30]
			(\x,{1.6*exp(-0.9*\x)}) -- (4,0) -- cycle;
		\draw[thick,navy] (-.3,0) -- (0,0) plot[domain=0:4,samples=40]
			(\x,{1.6*exp(-0.9*\x)});
		\node[below] at (1.5,0) {\small $x$};
		\node[teal] at (2.9,.65) {\small $e^{-\lambda x}$};
	\end{tikzpicture}
\end{center}

\emph{Uniform on $[\alpha,\beta]$}: the density is the constant
$1/(\beta-\alpha)$ on the interval, zero outside. Probability is length
divided by $\beta-\alpha$. Tiny example: uniform on $[0,10]$,
$P(2\le X\le5)=3/10$.

\say The height $1/(\beta-\alpha)$ is not a choice: it is what makes the
rectangle under the flat line have area 1, width times height. In the
left picture the shaded rectangle has the same height everywhere, so
every interval of the same length gets the same probability, wherever it
sits. ``Uniform'' is the continuous version of ``all outcomes equally
likely''.

\emph{Exponential with rate $\lambda>0$}: $f(x)=\lambda e^{-\lambda x}$ for
$x\ge0$. Then $F(x)=1-e^{-\lambda x}$, and the tail is
\[ P(X>x)=e^{-\lambda x}. \]
It models waiting times: the next order, a default. Tiny example:
$\lambda=1$, $P(X>2)=e^{-2}\approx0.135$.

\say Where does $F$ come from? Integrate the density:
$\int_0^x\lambda e^{-\lambda t}\,dt=\big[-e^{-\lambda t}\big]_0^x=1-e^{-\lambda x}$.
The tail, the shaded area in the right picture, is $1-F(x)=e^{-\lambda x}$:
the probability of waiting longer than $x$. Remember this tail; it is the
fastest route to most exponential questions. The rate $\lambda$ says how
often the events come: the larger $\lambda$, the faster the curve drops,
and the shorter the typical wait.

\subsection{Exercise~\exref{order-arrivals} --- Waiting for the next order, finance (10 min)}\slide{12--13}

\scriptexercise{order-arrivals}

\say This is an exponential with rate $\lambda=2$. Try (a)--(c), three
minutes.

(a) Click. $f\ge0$, and
$\int_0^\infty 2e^{-2x}\,dx=\big[-e^{-2x}\big]_0^\infty=0-(-1)=1$.

\say The antiderivative of $2e^{-2x}$ is $-e^{-2x}$: differentiate it,
and the chain rule brings down the factor $-2$. At infinity, $e^{-2x}$
goes to 0; at 0 it is 1. Hence $0-(-1)=1$.

Click. For $x\ge0$,
\[ F(x)=\int_0^x2e^{-2t}\,dt=\big[-e^{-2t}\big]_0^x=1-e^{-2x}, \]
and $F(x)=0$ for $x<0$. Click. The tail is $P(X>x)=e^{-2x}$.

\onslide the density $2e^{-2x}$, starting at height 2 and falling fast;
the area to the right of 1 is shaded, with an arrow labeled $P(X>1)$.
\begin{center}
	\begin{tikzpicture}[xscale=2,yscale=1.2]
		\fill[teal!30] (1,0) -- plot[domain=1:2.4,samples=40]
			(\x,{2*exp(-2*\x)}) -- (2.4,0) -- cycle;
		\draw[->] (-.2,0) -- (2.6,0) node[right] {\small $x$};
		\draw[->] (0,-.1) -- (0,2.3) node[above] {\small $f(x)$};
		\draw[navy,thick] plot[domain=0:2.4,samples=50]
			(\x,{2*exp(-2*\x)});
		\foreach \x in {0.5,1,1.5,2}
		\draw (\x,.04) -- (\x,-.04) node[below] {\scriptsize $\x$};
		\draw (.04,2) -- (-.04,2) node[left] {\scriptsize $2$};
		\draw[teal,->] (1.6,.7) node[above] {\small $P(X>1)$}
			-- (1.2,.12);
	\end{tikzpicture}
\end{center}

\say The density starts at 2, above 1: fine for a density. The shaded
tail to the right of 1 is thin: long waits are rare.

\begin{cue}
	Slide 13, one part per click.
\end{cue}

\ask How big do you expect the shaded area to be: closer to $0.5$ or to
$0.1$? \emph{(Closer to $0.1$: the curve is already low at $x=1$.)}

(b) $P(X>1)=e^{-2}\approx0.135$. About a 13.5\% chance that a whole minute
passes without an order.

(c) $X$ is continuous, so $P(X<0.5)=F(0.5)=1-e^{-1}\approx0.632$.

(d) \ask Which conditional probability is this? \emph{($P(X>1.5\mid
	X>0.5)$.)} ``More than $1.5$'' already contains ``more than $0.5$'', so
\[ P(X>1.5\mid X>0.5)=\frac{P(X>1.5)}{P(X>0.5)}=\frac{e^{-3}}{e^{-1}}
	=e^{-2}\approx0.135. \]
The definition of conditional probability from last week: the probability
of both events over the probability of the condition. ``Both'' is simply
$X>1.5$, since it implies $X>0.5$. Then the tail formula twice:
$P(X>1.5)=e^{-2\cdot1.5}=e^{-3}$ and $P(X>0.5)=e^{-2\cdot0.5}=e^{-1}$.

\ask Before the last click: compare with (b). What do you notice?

Click. \say The same as (b)! Having waited half a minute does not change
the chance of waiting one more minute. The exponential distribution has
no memory.

In general, $P(X>s+t\mid X>s)=e^{-\lambda(s+t)}/e^{-\lambda s}=e^{-\lambda
t}=P(X>t)$: the exponents subtract, and the past waiting time $s$
disappears. A quiet market in the last half minute tells us nothing about
how much longer the quiet will last.

\begin{cue}
	Common mistake in (d): computing $P(X>1.5)=e^{-3}$ and forgetting the
	condition. Another one: computing $P(0.5<X\le1.5)$, the probability of
	the interval, instead of the conditional probability.
\end{cue}

\begin{bg}{order arrivals and memorylessness}
	In market microstructure, orders arriving at an exchange are often
	modeled as a \emph{Poisson process}: the waiting times between orders
	are independent and exponential, with rate $\lambda$ orders per unit of
	time (here 2 per minute).
	Memorylessness says that the market ``does not remember'' how long it
	has been quiet. Real order flow comes in bursts (clustering), so the
	exponential model is a first approximation; models such as Hawkes
	processes add self-excitation.

	\terms order flow; arrival rate; Poisson process; memoryless.

	\sq{Is memorylessness realistic for defaults?}
	\ans For short horizons it is a common working assumption (constant
	hazard rate in credit models); over longer horizons the rate changes
	with the economy.
\end{bg}

\subsection{The distribution of a function of a random variable (6 min)}\slide{14}

\begin{cue}
	Clock check: this slide is planned for minute 51. If you start it after
	about minute 56, give only the two numbers in part (c) of
	Exercise~\exref{square-of-uniform}.
\end{cue}

\say We know the distribution of $X$. What is the distribution of
$Y=g(X)$, for example $X^2$, or the value $1000(1+X)$ of an investment
with return $X$? The method: always go through the distribution
function.

\onslide the method in three steps, and the example $Y=2X+3$:
\[ F_Y(y)=P(2X+3\le y)=P\Big(X\le \frac{y-3}{2}\Big)=\frac{y-3}{2},
	\qquad f_Y(y)=\frac12 . \]

\board
\begin{enumerate}[nosep]
	\item $F_Y(y)=P(Y\le y)=P(g(X)\le y)$;
	\item rewrite $g(X)\le y$ as an event about $X$, and use the
	      distribution of $X$;
	\item if you need the density: $f_Y(y)=F_Y'(y)$.
\end{enumerate}

Why through $F$, and not straight through the density? Because $F_Y(y)$
is a probability, $P(Y\le y)$, and probabilities do not care how we
describe an event: ``$Y\le y$'' and the equivalent statement about $X$ are
the same event, so they have the same probability. A density is not a
probability, and it does not survive such a rewriting.

Example from the slide: $X$ uniform on $[0,1]$, $Y=2X+3$. For
$3\le y\le5$:
\[ F_Y(y)=P(2X+3\le y)=P\Big(X\le\frac{y-3}{2}\Big)=\frac{y-3}{2}, \]
so $f_Y(y)=\tfrac12$: $Y$ is uniform on $[3,5]$. A linear function keeps
the shape; other functions change it, as we see now.

Step by step: subtract 3 and divide by 2 on both sides of the
inequality; for the uniform $X$, $P(X\le t)=t$. The uniform is stretched
to twice the length, so its height halves: $\tfrac12$.

\subsection{Exercise~\exref{square-of-uniform} --- The square of a uniform random variable (10 min)}\slide{15--16}

\scriptexercise{square-of-uniform}

\say Try (a) and (b), three minutes. Start with $P(X^2\le y)$.

Click. \say The method again: first the distribution function, then
differentiate.

(a) Click. $Y$ is between 0 and 1. For $0\le y\le1$, since $X\ge0$,
\[ F_Y(y)=P(X^2\le y)=P\big(X\le\sqrt y\big)=\sqrt y, \]
and $F_Y(y)=0$ for $y<0$, $F_Y(y)=1$ for $y>1$.

\say In general $X^2\le y$ means $-\sqrt y\le X\le\sqrt y$; here $X$ is
never negative, so only $X\le\sqrt y$ is left. And for the uniform $X$,
$P(X\le t)=t$, with $t=\sqrt y$. Check the ends: $F_Y(0)=0$ and
$F_Y(1)=1$, as it should be.

(b) Click. Differentiate: $f_Y(y)=\dfrac{1}{2\sqrt y}$ for $0<y\le1$, zero
otherwise.

\say $\sqrt y=y^{1/2}$, and its derivative is $\tfrac12y^{-1/2}$. This
density goes to infinity near 0, but its area stays 1:
$\int_0^1\frac{1}{2\sqrt y}\,dy=\big[\sqrt y\big]_0^1=1$. A density may be
unbounded; it only has to be nonnegative with total area 1.

\begin{cue}
	Slide 16.
\end{cue}

\onslide the density $f_Y(y)=1/(2\sqrt y)$ on $(0,1]$: very high near 0,
falling to $0.5$ at $y=1$. The area over $[0,0.25]$ is shaded teal
(labeled $0.5$), the area over $[0.81,1]$ gold (labeled $0.1$). A dashed
gray line at height 1 is the uniform density, for comparison.
\begin{center}
	\begin{tikzpicture}[xscale=5,yscale=.75]
		\fill[teal!30] (0,0) -- (0,{1/(2*sqrt(.016))}) -- plot[domain=.016:.25,samples=50]
			(\x,{1/(2*sqrt(\x))}) -- (.25,0) -- cycle;
		\fill[gold!50] (.81,0) -- plot[domain=.81:1,samples=20]
			(\x,{1/(2*sqrt(\x))}) -- (1,0) -- cycle;
		\draw[->] (0,0) -- (1.15,0) node[right] {\small $y$};
		\draw[->] (0,0) -- (0,4.3) node[above] {\small $f_Y(y)$};
		\draw[navy,thick] plot[domain=.016:1,samples=60]
			(\x,{1/(2*sqrt(\x))});
		\draw[navy,thick] (1,.5) -- (1,0);
		\foreach \x/\l in {.25/0.25,.81/0.81,1/1}
		\draw (\x,.08) -- (\x,-.08) node[below] {\scriptsize $\l$};
		\draw (.02,1) -- (-.02,1) node[left] {\scriptsize $1$};
		\draw[dashed,gray] (0,1) -- (1,1);
		\node[teal!60!black] at (.15,.45) {\scriptsize $0.5$};
		\node at (.9,.25) {\scriptsize $0.1$};
	\end{tikzpicture}
\end{center}

\ask Look at the two shaded areas. Which is bigger, and why, when the
gold interval is not much shorter than the teal one? \emph{(The teal one:
the density is much higher near 0.)}

(c) $P(Y\le0.25)=\sqrt{0.25}=0.5$, and $P(Y>0.81)=1-\sqrt{0.81}=0.1$.
Check: $Y\le0.25$ means $X\le0.5$; $Y>0.81$ means $X>0.9$.

\say Half of the probability lands in the first quarter of the interval.

(d) \ask Is $Y$ uniform? \emph{(No: the density $1/(2\sqrt y)$ is not
	constant.)} Click. Squaring pushes numbers in $(0,1)$ toward 0, for
example $0.5^2=0.25$, so small values of $Y$ are more likely.

\say In the picture, the curve is above the dashed uniform line for
$y<0.25$ and below it for $y>0.25$. Squaring does not create or destroy
probability; it moves it toward 0.

\begin{cue}
	Common mistake: $f_Y(y)=f_X(\sqrt y)=1$, plugging into the density of
	$X$; that would make $Y$ uniform. Always go through $F_Y$.
\end{cue}

\subsection{Summary (3 min)}\slide{17}

\onslide five bullet points with the formulas of this session; they are
the five lines you say now.

\say Today's first session in five lines: probability is the area under
the density; $F(x)=P(X\le x)$ and $f=F'$; uniform means length over total
length; the exponential has the tail $e^{-\lambda x}$ and no memory; and
for $Y=g(X)$, first $F_Y$, then differentiate. After the break: two random
variables at the same time.

\begin{cue}
	Use the reserve for questions, or let them start a self-study exercise
	on densities or on $Y=g(X)$ and walk around.
\end{cue}

%======================================================================
\section{Session 2: Joint distributions and independent random variables}
%======================================================================

\subsection{Session 2 title and recap of Session 1 (4 min)}\slide{18--19}

\begin{cue}
	Slide 18 (title), then slide 19.
\end{cue}

\onslide four bullet points, the summary of Session 1 again.

\say Welcome back. Before the break: densities, $F$, uniform and
exponential, and $Y=g(X)$. Now two random variables at the same time: how
do two stocks, two loans, two insurance claims behave \emph{together}?

Knowing each one alone is not enough. Two loans that each default with
probability 2\% can default together very rarely, or quite often, if a
recession hits both borrowers. To describe this we need the
\emph{joint} distribution: the probability of every combination.

\subsection{Joint distribution: the discrete case (6 min)}\slide{20}

\onslide three blocks of formulas, the joint pmf, the marginal pmfs
and the joint distribution function,
\[ p_X(x)=\sum_y p(x,y),\qquad p_Y(y)=\sum_x p(x,y),\qquad
	F(a,b)=P(X\le a,\,Y\le b), \]
and next to them a small table with the sums in the margins:
\begin{center}
	\renewcommand{\arraystretch}{1.25}
	\begin{tabular}{c|cc|c}
		$X\backslash Y$ & $0$ & $1$ & $p_X$\\\hline
		$0$ & $0.3$ & $0.2$ & \textcolor{teal}{$0.5$}\\
		$1$ & $0.1$ & $0.4$ & \textcolor{teal}{$0.5$}\\\hline
		$p_Y$ & \textcolor{teal}{$0.4$} & \textcolor{teal}{$0.6$} & $1$
	\end{tabular}\\[1mm]
	{\small the sums are written in the \emph{margins}: hence
		``marginal''; \ here $F(0,1)=0.3+0.2=0.5$}
\end{center}

\say If $X$ and $Y$ are discrete, their joint pmf is
\[ p(x,y)=P(X=x,\,Y=y). \]
It is a table: one cell for each pair of values, and all cells add up to
1.

\board Copy the small table of the slide.

\say The distribution of $X$ alone, the \emph{marginal} pmf $p_X$, is the
row sum: $p_X(x)=\sum_y p(x,y)$. The marginal of $Y$ is the column sum. We
write the sums in the margins of the table: hence the name.

Why the row sum? The event $X=0$ splits into disjoint pieces according to
the value of $Y$: $X=0$ and $Y=0$, or $X=0$ and $Y=1$. Disjoint pieces:
add their probabilities. That is the addition rule of last week, nothing
new.

\ask In the table, what is $P(X=0)$? \emph{($0.3+0.2=0.5$.)}

\ask And $P(Y=1)$? \emph{($0.2+0.4=0.6$, the second column.)}

\say Finally, the joint distribution function, for any $X$, $Y$:
$F(a,b)=P(X\le a,\,Y\le b)$. Here $F(0,1)=0.3+0.2=0.5$: the cells with
$X\le0$ and $Y\le1$.

To compute $F(a,b)$ from a table, add all cells that are ``up to $a$'' in
$X$ and ``up to $b$'' in $Y$.

\subsection{Exercise~\exref{two-stocks} --- Two stocks on the same day, finance (10 min)}\slide{21--22}

\scriptexercise{two-stocks}

\onslide the table under the statement (the margins stay empty until the
clicks of part (a)):
\begin{center}
	\small
	\begin{tabular}{c|ccc}
		       & $y=-1$ & $y=0$  & $y=1$  \\\hline
		$x=-1$ & $0.20$ & $0.10$ & $0.10$ \\
		$x=1$  & $0.05$ & $0.15$ & $0.40$
	\end{tabular}
\end{center}

\say Before you start, read the table aloud with me: the cell $0.40$ in
the bottom right is ``A rises and B rises''; the $0.20$ in the top left is
``both fall''. Try (a)--(c), three minutes.

(a) Click. \say Marginals: add up the rows and the columns, and write the
sums in the margins. Click. Row sums: $p_X(-1)=0.2+0.1+0.1=0.4$,
$p_X(1)=0.05+0.15+0.4=0.6$. Click. Column sums: $p_Y(-1)=0.25$,
$p_Y(0)=0.25$, $p_Y(1)=0.5$.

\onslide after the clicks of (a), the margins are filled in:
\begin{center}
	\renewcommand{\arraystretch}{1.25}
	\begin{tabular}{c|ccc|c}
		       & $y=-1$ & $y=0$  & $y=1$  & $p_X(x)$ \\
		\hline
		$x=-1$ & $0.20$ & $0.10$ & $0.10$ & \textcolor{teal}{$0.4$} \\
		$x=1$  & $0.05$ & $0.15$ & $0.40$ & \textcolor{teal}{$0.6$} \\
		\hline
		$p_Y(y)$ & \textcolor{teal}{$0.25$} & \textcolor{teal}{$0.25$}
		& \textcolor{teal}{$0.5$} & $1$
	\end{tabular}
\end{center}

\say Both margins add up to 1, a quick check of the arithmetic.

\begin{cue}
	Slide 22, one part per click.
\end{cue}

(b) \ask Which cells have $X=Y$? \emph{(Both fall, $(-1,-1)$, and both
	rise, $(1,1)$.)} So $P(X=Y)=0.20+0.40=0.60$.

(c) $X\le0$ means $X=-1$; $Y\le0$ means $Y=-1$ or $0$:
$F(0,0)=0.20+0.10=0.30$.

\ask Before (d): the probability that B rises is $0.5$. If you know that
A rose today, do you expect the chance for B to be higher, lower, or the
same? \emph{(Higher: the big number $0.40$ sits in the up--up cell.)}

(d) Conditional probability: the cell divided by the marginal,
\[ P(Y=1\mid X=1)=\frac{p(1,1)}{p_X(1)}=\frac{0.40}{0.60}=\frac23
	\approx0.667. \]
The definition from last week, $P(E\mid F)=P(E\cap F)/P(F)$: the
condition ``A rises'' keeps only the bottom row, which has total
probability $0.60$, and within it, the cell $0.40$.

Click. \say Compare with $P(Y=1)=0.5$. Knowing that A rises makes it more
likely that B rises too.

\begin{cue}
	Common mistake in (d): answering $p(1,1)=0.40$, without dividing by
	$p_X(1)=0.60$. In (c), $Y\le0$ includes $Y=0$.
\end{cue}

\begin{bg}{co-movement of stocks}
	Two stocks of the same market are exposed to common factors: the
	overall market, interest rates, the sector. So they tend to rise and
	fall on the same days. This co-movement is what correlation measures,
	and it is the reason diversification has limits: if all assets fall
	together, holding many of them does not help much. Here the up--up and
	down--down cells carry $0.60$ of the probability.

	\terms co-movement; common factor; market (systematic) risk;
	diversification.

	\sq{Can we read the co-movement from the marginals?}
	\ans No. The marginals say how each stock behaves alone; only the joint
	table says how they move together. Exercise~\exref{two-stocks-independence}
	shows this.
\end{bg}

\subsection{Joint distribution: the continuous case (6 min)}\slide{23}

\begin{cue}
	Clock check: this slide is planned for minute 20. If you start it after
	about minute 25, skip part (d) of Exercise~\exref{two-default-times};
	after about minute 30, also do part (d) of
	Exercise~\exref{joint-density-square} in one sentence.
\end{cue}

\say Now the same ideas for continuous random variables: the table
becomes a surface, sums become integrals.

\onslide three blocks, the joint density, the marginal densities and
$F(a,b)$; their formulas are displayed below. The picture: the unit square in the $(x,y)$-plane, lightly shaded, and in
it a darker triangle $A$, the part below the diagonal; ``the probability of
a region $A$ is the volume above it; do the inner integral first''.
\begin{center}
	\begin{tikzpicture}[scale=1.7]
		\fill[teal!20] (0,0) rectangle (1.4,1.4);
		\fill[teal!50] (0,0) -- (1.4,0) -- (1.4,1.4) -- cycle;
		\draw[->] (-.1,0) -- (1.7,0) node[right] {$x$};
		\draw[->] (0,-.1) -- (0,1.7) node[above] {$y$};
		\draw[thick,navy] (0,0) rectangle (1.4,1.4);
		\node[below] at (1.4,0) {\small 1};
		\node[left] at (0,1.4) {\small 1};
		\node at (1.0,.35) {\small $A$};
	\end{tikzpicture}
\end{center}

\say For two continuous random variables, a joint density $f(x,y)\ge0$
plays the role of the table. The probability of a region $A$ in the plane
is the double integral
\[ P\big((X,Y)\in A\big)=\iint_A f(x,y)\,dx\,dy, \]
the volume above $A$ under the surface, and the total is 1.

Picture the density as a tent or a roof over the square in the picture.
The probability that the point $(X,Y)$ lands in the dark triangle $A$ is
the volume of the part of the tent standing above the triangle.

\say Marginal densities: integrate the other variable out,
\[ f_X(x)=\int_{-\infty}^{\infty}f(x,y)\,dy,\qquad
	f_Y(y)=\int_{-\infty}^{\infty}f(x,y)\,dx. \]
It is the same as the row sum in the table, with an integral in place of
the sum. And
\[ F(a,b)=\int_{-\infty}^a\int_{-\infty}^b f(x,y)\,dy\,dx. \]

Tiny example: $f(x,y)=1$ on the unit square. Then
$P(X\le\tfrac12,\,Y\le\tfrac12)$ is the area of the small square,
$\tfrac14$. With a constant density 1, volume is just height 1 times area,
so probabilities are areas. The triangle $A$ of the picture, where
$X>Y$, has area $\tfrac12$, so $P(X>Y)=\tfrac12$. And the marginal is
$f_X(x)=\int_0^1 1\,dy=1$: $X$ alone is uniform on $[0,1]$.

\say For a double integral, do the inner integral first, with the other
variable held fixed.

\board $\displaystyle\int_0^1\!\!\int_0^1 f(x,y)\,dy\,dx$: inner over $y$
with $x$ fixed, then outer over $x$.

\subsection{Exercise~\exref{joint-density-square} --- A joint density on the unit square (12 min)}\slide{24--25}

\scriptexercise{joint-density-square}

\say Try (a) and (b), four minutes. Inner integral first, over $y$, with
$x$ fixed.

(a) Click. Inner: $\int_0^1(x+y)\,dy=\big[xy+\tfrac{y^2}{2}\big]_0^1=x+\tfrac12$.
Outer: $\int_0^1(x+\tfrac12)\,dx=\big[\tfrac{x^2}{2}+\tfrac x2\big]_0^1
=\tfrac12+\tfrac12=1$.

\say With $x$ held fixed, $x$ behaves like a number, so its integral in
$y$ is $xy$. Together with $f\ge0$ on the square, the total of 1 tells us
$f$ is a joint density.

\ask Before (b): with the small square being a quarter of the big one, do
you expect $\tfrac14$, more, or less? \emph{(Less: $x+y$ is smallest near
the origin, so the small square sits under the low part of the roof.)}

(b) Click. The same over the small square $[0,\tfrac12]\times[0,\tfrac12]$:
inner $\int_0^{1/2}(x+y)\,dy=\tfrac x2+\tfrac18$, outer
\[ \int_0^{1/2}\Big(\frac x2+\frac18\Big)dx
	=\Big[\frac{x^2}{4}+\frac{x}{8}\Big]_0^{1/2}=\frac1{16}+\frac1{16}=\frac18. \]

\say One eighth, half of $\tfrac14$: the guess was right. This number is
the value of the joint distribution function, $F(\tfrac12,\tfrac12)$.

\begin{cue}
	Slide 25.
\end{cue}

\onslide the unit square with the diagonal $y=x$ (dashed). The small
square $[0,\tfrac12]\times[0,\tfrac12]$ is outlined in gold, labeled (b);
the triangle below the diagonal is shaded teal, labeled (d).
\begin{center}
	\begin{tikzpicture}[scale=2.6]
		\fill[teal!30] (0,0) -- (1,0) -- (1,1) -- cycle;
		\fill[gold!55,opacity=.5] (0,0) rectangle (.5,.5);
		\draw[gold!80!black,thick] (0,0) rectangle (.5,.5);
		\draw[navy,thick] (0,0) rectangle (1,1);
		\draw[navy,dashed] (0,0) -- (1,1) node[above right] {\scriptsize $y=x$};
		\draw[->] (0,0) -- (1.2,0) node[right] {\small $x$};
		\draw[->] (0,0) -- (0,1.2) node[above] {\small $y$};
		\node[below] at (.5,0) {\scriptsize $\frac12$};
		\node[below] at (1,0) {\scriptsize $1$};
		\node[left] at (0,.5) {\scriptsize $\frac12$};
		\node[left] at (0,1) {\scriptsize $1$};
		\node at (.15,.37) {\scriptsize (b)};
		\node at (.75,.3) {\scriptsize (d)};
	\end{tikzpicture}
\end{center}

(c) \ask Haven't we done this already? \emph{(Yes: the inner integral in
	(a).)} $f_X(x)=x+\tfrac12$ for $0\le x\le1$, zero otherwise.

\say Integrating out $y$ is exactly what the inner integral of (a) did.
A quick check: $\int_0^1(x+\tfrac12)\,dx=1$, the outer integral of (a).
By the same computation, $f_Y(y)=y+\tfrac12$. Keep these two marginals in
mind; we will need them in a few minutes for independence.

(d) Click. \ask Do we need a double integral over the triangle? \emph{(No.)}
Swapping $x$ and $y$ does not change $f$, so $P(X>Y)=P(Y>X)$; and
$P(X=Y)=0$, a line has no area. The two add up to 1, so
$P(X>Y)=\tfrac12$.

\say In the picture: the teal triangle below the diagonal is $X>Y$, its
mirror image above the diagonal is $Y>X$. The roof $x+y$ looks the same
from both sides of the diagonal, so the two volumes are equal. The
diagonal itself is a line: no area, no volume.

\begin{cue}
	Common mistake in (b): taking the area $\tfrac14$ of the small square;
	that works only for a constant density. If someone integrates over the
	triangle in (d) ($y$ from $0$ to $x$), let them: it also gives
	$\tfrac12$, since $\int_0^1\int_0^x(x+y)\,dy\,dx=\int_0^1\tfrac32x^2\,dx
		=\tfrac12$.
\end{cue}

\subsection{Independent random variables (6 min)}\slide{26}

\begin{cue}
	Clock check: this slide is planned for minute 38. If you start it after
	about minute 43, give part (c) of
	Exercise~\exref{two-stocks-independence} as one number and one
	sentence, and skip part (d) of Exercise~\exref{two-default-times} if
	you have not decided so yet.
\end{cue}

\say Last week, two \emph{events} were independent if the probability of
both is the product. Now we lift this to random variables.

\onslide the definition (below) and two blocks with the practical tests:
\[ \text{discrete: } p(x,y)=p_X(x)\,p_Y(y)\ \text{in every cell};\qquad
	\text{continuous: } f(x,y)=f_X(x)\,f_Y(y). \]

\say $X$ and $Y$ are independent if every event about $X$ is independent
of every event about $Y$:
\[ P(X\in A,\,Y\in B)=P(X\in A)\,P(Y\in B). \]
Equivalently, $F(a,b)=F_X(a)F_Y(b)$ for all $a,b$.

In practice: discrete, the product rule must hold in \emph{every} cell;
one cell that fails, and they are not independent. Continuous, the joint
density is the product of the marginals; on a rectangle, $f$ splits into a function of $x$
times a function of $y$.

Tiny example, discrete: the small table of slide 20. The top left cell is
$p(0,0)=0.3$, but $p_X(0)\,p_Y(0)=0.5\cdot0.4=0.2$. One cell fails: not
independent.

Tiny example, continuous: $f(x,y)=4xy$ on the unit square is
$2x\cdot2y$: independent; its marginals are $2x$ and $2y$.
$f(x,y)=x+y$ of the last exercise does not split: not independent. Its
marginals were $x+\tfrac12$ and $y+\tfrac12$; at $(0,0)$ their product
is $\tfrac14$, but $f(0,0)=0$.

\say In words: knowing the value of $X$ does not change the distribution
of $Y$.

\subsection{Exercise~\exref{two-stocks-independence} --- Do the two stocks move independently?, finance (8 min)}\slide{27--28}

\scriptexercise{two-stocks-independence}

\say The same two stocks as before. Try (a), two minutes: compare one cell
with the product of its margins.

\ask Which cell would you test first? \emph{(A cell with a big or a small
number, for example $0.20$ or $0.40$: those are the most likely to fail.)}

(a) Click. $p(-1,-1)=0.20$, but $p_X(-1)p_Y(-1)=0.4\cdot0.25=0.10$. Not
independent. Click. \say One bad cell is enough for ``no''; for ``yes'' all
cells must match. The cell $(-1,0)$ happens to match, $0.10=0.4\cdot0.25$,
and proves nothing.

\say The cell ``both fall'' has twice the probability that independent
stocks would give it: the two stocks fall together more often than
chance alone would explain.

\begin{cue}
	Slide 28.
\end{cue}

(b) \ask How do we fill the table of independent stocks? \emph{(Each cell
	is row sum times column sum.)} Click. For example $0.4\cdot0.5=0.20$ and
$0.6\cdot0.25=0.15$.

\onslide the table of independent stocks with the same margins (the cells
appear at the click):
\begin{center}
	\renewcommand{\arraystretch}{1.25}
	\begin{tabular}{c|ccc|c}
		       & $y=-1$ & $y=0$  & $y=1$  & $p_X(x)$ \\
		\hline
		$x=-1$ & $0.10$ & $0.10$ & $0.20$ & $0.4$ \\
		$x=1$  & $0.15$ & $0.15$ & $0.30$ & $0.6$ \\
		\hline
		$p_Y(y)$ & $0.25$ & $0.25$ & $0.5$ & $1$
	\end{tabular}
\end{center}

\say Check the new table: each row still adds up to its margin,
$0.10+0.10+0.20=0.4$ and $0.15+0.15+0.30=0.6$, and each column too. So the
two stocks behave exactly as before \emph{one at a time}; only the way
they move together has changed.

\ask Before (c): for these independent stocks, will $P(X=Y)$ be bigger or
smaller than the $0.60$ of the real stocks? \emph{(Smaller.)}

(c) Click. $P(X=Y)=0.10+0.30=0.40$, against $0.60$ for the real stocks.
Click. \say The real stocks move in the same direction more often than
independent ones would: they tend to move together.

\say And this is the lesson of the whole exercise: the two tables have
the same margins, yet $0.60$ against $0.40$. The marginals alone cannot
tell how two random variables move together; for that we need the joint
distribution.

\begin{cue}
	Common mistake in (a): testing only the cell $(-1,0)$, finding a match,
	and concluding ``independent''.
\end{cue}

\begin{bg}{correlation and diversification}
	The table of (b) is what the stocks would look like with the same
	individual behavior but no common driver. The gap between $0.60$ and
	$0.40$ is the dependence. In portfolio theory, dependence (measured by
	correlation) decides how much risk diversification removes: with
	independent assets, risk shrinks as more are added; with assets that
	move together, a part of the risk always remains.

	\terms independence; correlation; diversification.

	\sq{Are stock returns ever independent?}
	\ans Almost never exactly: most stocks share exposure to the market.
	Independence is a modeling simplification.
\end{bg}

\subsection{Exercise~\exref{two-default-times} --- Two borrowers, independent defaults, finance (10 min)}\slide{29--30}

\scriptexercise{two-default-times}

\say Try (a)--(c), three minutes. Remember the tail of the exponential.

(a) Click. $P(T>t)=e^{-\lambda t}$. Click. So
$P(T_A>10)=e^{-0.2}\approx0.819$ and $P(T_B>10)=e^{-0.3}\approx0.741$.

\say The exponent is rate times time: $0.02\cdot10=0.2$ and
$0.03\cdot10=0.3$. Company B, with the higher rate, is the riskier one.

(b) Click. Independent, so the probabilities multiply:
\[ P(T_A>10,\,T_B>10)=e^{-0.2}e^{-0.3}=e^{-0.5}\approx0.607. \]

\begin{cue}
	Slide 30.
\end{cue}

\onslide the $(t_A,t_B)$-plane, split by dashed lines at $t_A=10$ and
$t_B=10$. The upper right region ($T_A>10$ and $T_B>10$) is shaded teal,
``both survive''; from the third click, the lower left square ($T_A\le10$
and $T_B\le10$) is shaded gold, ``both default''.
\begin{center}
	\begin{tikzpicture}[scale=.6]
		\fill[teal!20] (2.5,2.5) rectangle (6,6);
		\fill[gold!50] (0,0) rectangle (2.5,2.5);
		\draw[->] (0,0) -- (6.3,0) node[below left,font=\small] {$t_A$};
		\draw[->] (0,0) -- (0,6.3) node[below left,font=\small] {$t_B$};
		\draw[dashed] (2.5,0) -- (2.5,6) (0,2.5) -- (6,2.5);
		\node[below,font=\small] at (2.5,0) {$10$};
		\node[left,font=\small] at (0,2.5) {$10$};
		\node[font=\small,align=center] at (4.25,4.25) {both\\survive};
		\node[font=\small,align=center] at (1.25,1.25) {both\\def.};
	\end{tikzpicture}
\end{center}

(c) \ask What is the complement of ``at least one defaults''? \emph{(Both
	survive.)} So $1-e^{-0.5}\approx0.393$.

\say In the picture: ``at least one defaults'' is everything outside the
teal region, the three other pieces of the plane together. Its
probability is one minus the teal one.

(d) Click. $P(T_A\le10)\approx0.181$, $P(T_B\le10)\approx0.259$. Click. By
independence the product, $\approx0.047$.

\say The gold square in the picture. Each default probability is one
minus the survival probability: $1-0.819$ and $1-0.741$. A check that ties
(c) and (d) together: by inclusion--exclusion,
$P(\text{at least one})=0.181+0.259-0.047\approx0.393$, the answer of (c).

\ask Is $0.047$, under 5\%, realistic for two borrowers in the same
economy? \emph{(Probably too low: in a recession, both get into trouble at
the same time.)}

Click. \say In a crisis, defaults come together, so independence underestimates
(d).

\begin{cue}
	Common mistake in (c): $P(T_A\le10)+P(T_B\le10)\approx0.440$, which
	counts ``both default'' twice.
\end{cue}

\begin{bg}{default times and default correlation}
	In reduced-form credit models the default time of a company is
	exponential with a \emph{hazard rate} (intensity) $\lambda$: the rate
	$0.02$ means roughly a 2\% chance of default per year. Survival
	probability is $e^{-\lambda t}$. The probability that \emph{both}
	default is what matters for a lender with exposure to both, and for
	products such as first-to-default baskets and CDO tranches. Assuming
	independence understates joint defaults; in 2007--2008 this was one
	reason why such products lost far more than models predicted.

	\terms hazard rate; survival probability; default correlation;
	first-to-default basket.

	\sq{How do practitioners model dependent defaults?}
	\ans With common factors or copulas: the marginals stay exponential,
	but the joint distribution makes defaults cluster.
\end{bg}

\subsection{Summary (3 min)}\slide{31}

\onslide five bullet points with the formulas of this session; they are
the lines you say now.

\say Session 2 in a few lines: a discrete joint distribution is a table,
and the marginals are its row and column sums; a continuous one is a joint
density, probabilities are double integrals, and the marginals come from
integrating the other variable out; $F(a,b)=P(X\le a,\,Y\le b)$; and
independent means the joint pmf or density is the product of the
marginals, so probabilities of ``and'' multiply. The marginals alone do
not tell how $X$ and $Y$ move together.

\say Every exercise of today
is in the study guide on the course website, with a step-by-step
solution, and the Exercises page there has more problems on every topic
for self-study. Thank you, and see you next week.

\begin{cue}
	Use the reserve for questions, or let them start a self-study exercise
	on independence and walk around.
\end{cue}

\end{document}
